\section{敏感性、稳健性与一致性检验}
\subsection{时间标签敏感性}
将问题一的144点净负荷与电价整体平移$-10$、0、$+10$分钟并重新优化，费用分别为35,129.73、35,126.95和35,126.85元；相对主口径最大偏差仅2.78元，说明总费用对这一行序歧义不敏感。指定时段的解释仍应遵循官方模板行序。
\subsection{可行性与数值残差}
求解前以零储能、直接购电构造可行基线；求解后独立回代供需、SOC、容量边界及费用分量。问题一与334日滚动任务共1,337项检查全部通过。最小供给松弛为$-1.137\times10^{-13}$ kWh，最大SOC递推残差为$1.819\times10^{-12}$ kWh，最大同时充放电量为0；所有SOC均位于$[1200,10800]$ kWh。上述误差处于浮点求解容差内。
\subsection{结果稳健性边界}
问题二的结论同时体现在费用下降22.96\%和紧急购电下降42.15\%，方向一致。问题三及问题4-3则表现为紧急购电下降、费用上升，说明其结论依赖调整费结构。当前实现仅比较“全部更新”与“完全不调整”两个端点，尚未穷举0/6、0/12、0/18等发布组合；完整场景CVaR、不同调整结算解释和充放电效率口径应作为后续扩展，而不能据现有结果虚构敏感性数值。
